\chapter{Testicular infection}
\index{Testicular infection}

\medskip
\noindent\textit{Educational reference; always seek professional advice for individual care.}

\section*{Definition and Overview}
Testicular infection refers to inflammation and infection of the male reproductive organs in the scrotum. The condition is usually one of three related problems, which commonly occur together. Orchitis is inflammation of the testis itself, epididymitis is inflammation of the epididymis, the coiled tube at the back of the testis that stores and carries sperm, and epididymo-orchitis is inflammation of both together, which is the most common form. Most cases in younger, sexually active men are caused by sexually transmitted bacteria such as chlamydia and gonorrhoea, while in older men the cause is more often a urinary tract infection spreading from the prostate or bladder. Viral orchitis, most famously from the mumps virus, also occurs and causes a very painful, swollen testis. A painful, swollen scrotum is a symptom that must always be taken seriously, because the most urgent differential diagnosis, testicular torsion, can destroy a testis within hours if not treated. This chapter explains the causes, symptoms, diagnosis, and treatment of testicular infection, and clarifies the difference between it and other causes of acute scrotal pain.

\section*{Epidemiology}
Epididymitis and epididymo-orchitis are among the most common causes of acute scrotal pain in adult men, accounting for the majority of scrotal presentations to general practitioners and emergency departments. The peak age for sexually transmitted epididymitis is between 18 and 35 years, matching the age range in which chlamydia and gonorrhoea are most common. In men over 40, infection more often follows a urinary tract infection or prostatitis, and after prostate surgery. Viral orchitis is most commonly seen in unvaccinated adolescents and young adults, since the introduction of the measles-mumps-rubella (MMR) vaccine has made mumps orchitis much rarer in developed countries; it typically affects about one in four males who develop mumps after puberty. Bacterial orchitis without epididymitis is uncommon, because the testis itself is relatively protected, and infection usually reaches it through the epididymis. Overall, epididymitis and epididymo-orchitis affect men of all ages and are more common than is generally appreciated.

\section*{Aetiology and Pathophysiology}
Infection reaches the epididymis and testis by one of two main routes.

\begin{itemize}
    \item \textbf{Ascending infection from the urinary tract:} bacteria travel from the bladder, prostate, or urethra up the vas deferens into the epididymis. This is the mechanism in most cases
    \item \textbf{Blood-borne spread:} in viral orchitis such as mumps, the virus reaches the testis through the bloodstream during a generalised infection
\end{itemize}

The organisms involved differ with age and sexual activity:
\begin{itemize}
    \item \textbf{In sexually active men under 35:} \textit{Chlamydia trachomatis} and \textit{Neisseria gonorrhoeae} are the commonest causes, usually as part of a sexually transmitted infection
    \item \textbf{In older men:} coliform bacteria such as \textit{Escherichia coli} from the bowel, spreading from an infected prostate or bladder, are typical
    \item \textbf{In men who practise anal sex:} bowel organisms may also cause infection
    \item \textbf{Viral causes:} the mumps virus is the classic cause of viral orchitis; influenza and other viruses can occasionally do the same
    \item \textbf{Rarer causes:} tuberculosis, fungal infection, and, very rarely, amiodarone, a heart medication that can cause a drug-related epididymitis
\end{itemize}

Once bacteria reach the epididymis, they provoke intense inflammation, swelling, and pain, and the infection can then extend into the testis itself. A coexisting structural problem, such as an enlarged prostate or a urinary catheter, increases the risk by making the urinary tract more hospitable to bacteria.

\section*{Clinical Features}
Testicular infection typically comes on over one to three days and is most commonly one-sided. The typical features are:

\begin{itemize}
    \item \textbf{Gradually worsening pain and swelling} in one testicle, which can be mild or severe, and may spread to the groin or lower abdomen
    \item \textbf{A swollen, tender, warm, red scrotum:} the affected testis is often markedly enlarged and painful to touch
    \item \textbf{Urinary symptoms:} pain or burning when passing urine, frequency, urgency, or a urethral discharge, particularly with sexually transmitted causes
    \item \textbf{Fever and a general feeling of being unwell:} more common in bacterial infection
    \item \textbf{A visible discharge from the penis} in sexually transmitted infection
    \item \textbf{In mumps orchitis:} fever, swollen cheeks from parotitis, and a tender, swollen testis appearing days later, sometimes affecting both sides
    \item \textbf{Relief with elevation:} unlike testicular torsion, the pain of epididymitis is often eased slightly by lifting the scrotum
\end{itemize}

It is essential to remember that a sudden, very severe, one-sided scrotal pain in a child or young man is testicular torsion until proven otherwise, and requires emergency assessment regardless of how it resembles infection.

\section*{Differential Diagnosis}
An acutely painful, swollen scrotum is a genuine emergency scenario, and the priority is to exclude conditions that threaten the testis.

\begin{itemize}
    \item \textbf{Testicular torsion:} twisting of the spermatic cord cuts off the blood supply and can destroy the testis within four to six hours. It typically causes sudden, severe pain, often with nausea and vomiting, and is more common in boys and young men. It is a surgical emergency
    \item \textbf{Torsion of the testicular appendage:} a small twist of the vestigial remnant of tissue on the testis, causing localised pain, most common in boys
    \item \textbf{Incarcerated inguinal hernia:} a loop of bowel trapped in the groin, causing a lump and pain
    \item \textbf{Testicular cancer:} usually a painless lump, but a tumour can sometimes present with pain or following trauma
    \item \textbf{Trauma:} direct injury to the scrotum, which can also rupture the testis
    \item \textbf{Fournier's gangrene:} a rare but life-threatening, rapidly spreading infection of the scrotal skin and underlying tissues, usually in older men with diabetes
\end{itemize}

Because several of these require urgent surgery, imaging is often performed before any treatment decision.

\section*{When to Seek Medical Care}
Any man or boy with new scrotal pain or swelling should be assessed promptly. This is a case where it is always better to seek care early, because the serious alternative diagnoses are time-critical.

\textbf{Seek emergency care immediately, and contact local emergency services if appropriate, for:}
\begin{itemize}
    \item Sudden, severe, one-sided scrotal pain, especially in a child or adolescent
    \item Pain accompanied by nausea, vomiting, or collapse
    \item A high testicle or a "bell-clapper" appearance, or pain that started abruptly
    \item Rapidly spreading redness, fever, and feeling very unwell, which may indicate Fournier's gangrene
\end{itemize}

Non-urgent medical review is appropriate within a day or two for milder, gradually developing pain and swelling, urinary symptoms, or a urethral discharge. However, if there is any doubt, an emergency assessment is the safest course.

\section*{Diagnosis and Investigations}
The assessment of a painful scrotum begins with a careful history and examination, and the priority is to distinguish torsion from infection.

\begin{itemize}
    \item \textbf{Examination:} the clinician assesses the position and tenderness of the testis, the scrotal skin, and whether the pain eases when the scrotum is lifted (the Prehn sign), which suggests epididymitis
    \item \textbf{Urine tests:} a midstream urine sample is tested for white cells and bacteria, and is cultured in older men
    \item \textbf{Swabs:} urethral swabs or first-pass urine samples are tested for chlamydia and gonorrhoea in sexually active men under 35
    \item \textbf{Ultrasound with Doppler:} scans the blood flow to the testis; increased flow suggests infection and reduced or absent flow suggests torsion. This is the key test when there is doubt
    \item \textbf{Blood tests:} markers of inflammation such as the white cell count and CRP are raised in infection
    \item \textbf{Mumps serology:} blood tests for mumps antibodies may be arranged when viral orchitis is suspected
    \item \textbf{If torsion is strongly suspected:} urgent surgical exploration is performed rather than waiting for tests, because every hour matters
\end{itemize}

\section*{Management}
Treatment depends on the cause and severity, but early treatment speeds recovery and reduces complications.

\begin{itemize}
    \item \textbf{Antibiotics:} for bacterial epididymitis, a course of antibiotics is given. In sexually active young men this covers chlamydia and gonorrhoea (for example, doxycycline plus ceftriaxone), and in older men it covers bowel bacteria. Sexual partners are treated too, to prevent reinfection
    \item \textbf{Pain relief:} simple analgesia such as paracetamol and ibuprofen
    \item \textbf{Scrotal support and elevation:} wearing supportive underwear, and resting, ease the pain
    \item \textbf{Cold packs:} applied to the scrotum for short periods can reduce swelling
    \item \textbf{Avoidance of heavy lifting and vigorous activity} until symptoms settle
    \item \textbf{Mumps orchitis:} there is no specific antiviral treatment; management is supportive with rest, fluids, analgesia, and cold packs until the infection resolves
    \item \textbf{Treating the underlying cause:} if an enlarged prostate or urinary retention is driving recurrent infection, this is addressed separately
    \item \textbf{Referral:} complicated, recurrent, or severe cases may be referred to a urologist
    \item \textbf{Follow-up:} a review is arranged to confirm the swelling has resolved, and an ultrasound may be repeated if there is any concern
\end{itemize}

Most men improve within a few days of starting treatment, though tenderness can persist for some weeks.

\section*{Complications}
Most cases of testicular infection resolve completely, but complications can occur, particularly when treatment is delayed or infection is severe.

\begin{itemize}
    \item \textbf{Abscess formation:} a pocket of pus in the scrotum requiring drainage
    \item \textbf{Chronic epididymitis:} persistent pain and discomfort lasting months, which can be difficult to treat
    \item \textbf{Reduced fertility:} severe or bilateral orchitis can damage the seminiferous tubules and reduce sperm production; this is most significant after bilateral mumps orchitis
    \item \textbf{Testicular atrophy:} shrinkage of the testis after severe inflammation
    \item \textbf{Spread of infection:} ascending infection can reach the prostate or, rarely, cause systemic sepsis
    \item \textbf{Fournier's gangrene:} a rare but life-threatening necrotising infection that requires emergency surgery
    \item \textbf{Emotional and psychological effects:} pain, fertility concerns, and the association with STIs can cause distress, which is a legitimate part of care
\end{itemize}

\section*{Prognosis and Outcomes}
With early and appropriate treatment, the outlook for testicular infection is very good. Most men are significantly better within a few days and fully recovered within several weeks. Fertility is normally preserved after a single, well-treated episode in one testis, because the remaining healthy tissue compensates. The main adverse outcomes arise when torsion is missed (which is why the diagnosis is taken so seriously) or when infection is severe, recurrent, or bilateral. Chronic epididymitis develops in a minority and can cause ongoing discomfort, but it is manageable with a urologist's help. Mumps orchitis, when it affects both testes after puberty, carries a higher risk of reduced fertility. Overall, the message is reassuring: prompt assessment and treatment of a painful scrotum almost always leads to full recovery.

\section*{Prevention and Self-Care}
\begin{itemize}
    \item Practise safer sex using condoms to reduce the risk of chlamydia and gonorrhoea, the leading causes in young men
    \item If diagnosed with an STI, complete the full antibiotic course and inform partners so they can be treated
    \item Treat urinary tract infections promptly, and drink plenty of water to keep the urinary tract healthy
    \item Ensure MMR vaccination is up to date to prevent mumps and its complications
    \item Support the scrotum with supportive underwear during and after an episode
    \item Rest and avoid heavy lifting and vigorous exercise while symptoms settle
    \item Use cold packs and simple analgesia for pain, and seek review if symptoms worsen or persist
    \item Any new scrotal lump or persistent swelling, even if painless, should be assessed promptly to exclude testicular cancer
\end{itemize}

\section*{Sources and Further Reading}
The following organisations provide reliable information about testicular infection, acute scrotal pain, and testicular health.

\begin{itemize}
    \item NHS. \textit{Epididymitis and orchitis: overview and treatment.} https://www.nhs.uk/
    \item Mayo Clinic. \textit{Epididymitis: symptoms, causes, and treatment.} https://www.mayoclinic.org/
    \item MSD Manual. \textit{Epididymitis and orchitis: consumer and professional overview.} https://www.msdmanuals.com/
    \item MedlinePlus. \textit{Orchitis and epididymitis health information.} https://medlineplus.gov/
    \item Australian Department of Health and Aged Care. \textit{Sexually transmitted infections information.} https://www.health.gov.au/
    \item Testicular Cancer Australia. \textit{Testicular health and lump awareness.} https://www.testicularcancer.org.au/
\end{itemize}
